\documentclass[10pt,a4paper]{amsart}
\usepackage{amsmath,amssymb,amsthm,graphicx}
% Retained minimal preamble for the Stacks Project chapter varieties.tex.
% Verbatim excerpt of the chapter preamble: Xy-pic setup (native preamble lines 14-18), the native theorem
% declarations with the mathematical macros (native preamble lines 154-217), and the project own no-op
% definitions of the reference/slogan/history display environments (native wrapper.tex lines 1-3) that the
% counted statement uses for its slogan. Copied mechanically; nothing added, removed or redefined.
% For commutative diagrams we use Xy-pic
\usepackage[all]{xy}

% We use 2cell for 2-commutative diagrams.
\xyoption{2cell}
% Theorem environments.
%
\theoremstyle{plain}
\newtheorem{theorem}[subsection]{Theorem}
\newtheorem{proposition}[subsection]{Proposition}
\newtheorem{lemma}[subsection]{Lemma}

\theoremstyle{definition}
\newtheorem{definition}[subsection]{Definition}
\newtheorem{example}[subsection]{Example}
\newtheorem{exercise}[subsection]{Exercise}
\newtheorem{situation}[subsection]{Situation}

\theoremstyle{remark}
\newtheorem{remark}[subsection]{Remark}
\newtheorem{remarks}[subsection]{Remarks}

\numberwithin{equation}{subsection}

% Macros
%
\def\lim{\mathop{\mathrm{lim}}\nolimits}
\def\colim{\mathop{\mathrm{colim}}\nolimits}
\def\Spec{\mathop{\mathrm{Spec}}}
\def\Hom{\mathop{\mathrm{Hom}}\nolimits}
\def\Ext{\mathop{\mathrm{Ext}}\nolimits}
\def\SheafHom{\mathop{\mathcal{H}\!\mathit{om}}\nolimits}
\def\SheafExt{\mathop{\mathcal{E}\!\mathit{xt}}\nolimits}
\def\Sch{\mathit{Sch}}
\def\Mor{\mathop{\mathrm{Mor}}\nolimits}
\def\Ob{\mathop{\mathrm{Ob}}\nolimits}
\def\Sh{\mathop{\mathit{Sh}}\nolimits}
\def\NL{\mathop{N\!L}\nolimits}
\def\CH{\mathop{\mathrm{CH}}\nolimits}
\def\proetale{{pro\text{-}\acute{e}tale}}
\def\etale{{\acute{e}tale}}
\def\QCoh{\mathit{QCoh}}
\def\Ker{\mathop{\mathrm{Ker}}}
\def\Im{\mathop{\mathrm{Im}}}
\def\Coker{\mathop{\mathrm{Coker}}}
\def\Coim{\mathop{\mathrm{Coim}}}

% Boxtimes
%
\DeclareMathSymbol{\boxtimes}{\mathbin}{AMSa}{"02}

%
% Macros for moduli stacks/spaces
%
\def\QCohstack{\mathcal{QC}\!\mathit{oh}}
\def\Cohstack{\mathcal{C}\!\mathit{oh}}
\def\Spacesstack{\mathcal{S}\!\mathit{paces}}
\def\Quotfunctor{\mathrm{Quot}}
\def\Hilbfunctor{\mathrm{Hilb}}
\def\Curvesstack{\mathcal{C}\!\mathit{urves}}
\def\Polarizedstack{\mathcal{P}\!\mathit{olarized}}
\def\Complexesstack{\mathcal{C}\!\mathit{omplexes}}
% \Pic is the operator that assigns to X its picard group, usage \Pic(X)
% \Picardstack_{X/B} denotes the Picard stack of X over B
% \Picardfunctor_{X/B} denotes the Picard functor of X over B
\def\Pic{\mathop{\mathrm{Pic}}\nolimits}
\def\Picardstack{\mathcal{P}\!\mathit{ic}}
\def\Picardfunctor{\mathrm{Pic}}
\def\Deformationcategory{\mathcal{D}\!\mathit{ef}}
\newenvironment{reference}{}{}
\newenvironment{slogan}{}{}
\newenvironment{history}{}{}
\usepackage{cleveref}
\crefname{originallemma}{original lemma}{original lemmas}
\Crefname{originallemma}{Original lemma}{Original lemmas}
\makeatletter
\newlabel{algebra-lemma-geometrically-reduced-finite-purely-inseparable-extension}{{030V}{source}{}{sourcepointer-1}{}}
\newlabel{algebra-lemma-geometrically-reduced-finite-purely-inseparable-extension@cref}{{[originallemma][0][]030V}{[1][0][]source}{}{}{}}
\makeatother
\crefname{originallemma}{original lemma}{original lemmas}
\Crefname{originallemma}{Original lemma}{Original lemmas}
\makeatletter
\newlabel{algebra-lemma-p-ring-map}{{0BRA}{source}{}{sourcepointer-2}{}}
\newlabel{algebra-lemma-p-ring-map@cref}{{[originallemma][0][]0BRA}{[1][0][]source}{}{}{}}
\makeatother
\crefname{originallemma}{original lemma}{original lemmas}
\Crefname{originallemma}{Original lemma}{Original lemmas}
\makeatletter
\newlabel{lemma-geometrically-reduced-at-point}{{035W}{source}{}{sourcepointer-3}{}}
\newlabel{lemma-geometrically-reduced-at-point@cref}{{[originallemma][0][]035W}{[1][0][]source}{}{}{}}
\makeatother
\title{Geometric reducedness at a point tested over purely inseparable extensions and over the perfect closure}
\author{The Stacks Project Authors}
\date{stacks-project:varieties.tex:local-snapshot-20261011; GNU FDL 1.2 or later}
\begin{document}
\maketitle
\noindent\textit{Editorial note.} All delivered mathematical text is the author's own, copied from the pinned Stacks Project chapter file varieties.tex; no statement or proof step was rewritten, reordered or removed.\par\medskip
\noindent\textit{Editorial note.} The proof is complete as written and covers all four conditions through the two cited criteria; no step is referred to the reader.\par\medskip
\noindent\textit{Editorial note.} The counted spans carry 3 navigation stubs and no bibliography citation.\par\medskip
\noindent\textit{Editorial note.} The preamble is the chapter's retained minimal preamble (sources/stacks-varieties/retained-preamble.tex lines 1-77), a verbatim excerpt of the native preamble.\par\medskip
\noindent\textit{Editorial note.} The label lemma-check-only-finite-inseparable-extensions is defined only in varieties.tex, so no cross-chapter restatement had to be weighed for this card.\par\medskip
\noindent\textit{Allowed standard prerequisites (editorial).}\begin{itemize}
\item A purely inseparable base change induces a homeomorphism and identifies the local ring with the tensor product. Algebra, Lemma algebra-lemma-p-ring-map, invoked at proof lines 479-484. Supplies the uniqueness of the lifted point and the identification of the local ring.
\item Pointwise criterion for geometric reducedness of a scheme at a point. Same chapter, Lemma lemma-geometrically-reduced-at-point, invoked at proof lines 485-486. Supplies the passage between the scheme statement and the local rings.
\item Geometric reducedness tested over finite purely inseparable extensions. Algebra, Lemma algebra-lemma-geometrically-reduced-finite-purely-inseparable-extension, invoked at proof lines 486-487. Supplies the algebraic equivalence.
\end{itemize}\medskip
\ifdefined\hypertarget\hypertarget{sourcepointer-1}{}\fi
\noindent\textit{Source reference.} Original lemma 030V: sources/stacks-chow/tags.txt, source lines 3962--3962. Pointer outside the counted span, invoked once for context; the load-bearing argument of the counted statement is the author's proof reproduced in full in the proof span.\par\medskip
\ifdefined\hypertarget\hypertarget{sourcepointer-2}{}\fi
\noindent\textit{Source reference.} Original lemma 0BRA: sources/stacks-chow/tags.txt, source lines 14763--14763. Pointer outside the counted span, invoked once for context; the load-bearing argument of the counted statement is the author's proof reproduced in full in the proof span.\par\medskip
\ifdefined\hypertarget\hypertarget{sourcepointer-3}{}\fi
\noindent\textit{Source reference.} Original lemma 035W: sources/stacks-varieties/pinned-original-source.tex, source lines 354--354. Pointer outside the counted span, invoked once for context; the load-bearing argument of the counted statement is the author's proof reproduced in full in the proof span.\par\medskip
\begin{lemma}
\label{lemma-check-only-finite-inseparable-extensions}
Let $k$ be a field of characteristic $p > 0$. Let $X$ be a scheme over $k$.
Let $x \in X$. The following are equivalent
\begin{enumerate}
\item $X$ is geometrically reduced at $x$,
\item $\mathcal{O}_{X_{k'}, x'}$ is reduced for every
finite purely inseparable field extension $k'$ of $k$ and
$x' \in X_{k'}$ the unique point lying over $x$,
\item $\mathcal{O}_{X_{k^{1/p}}, x'}$ is reduced for
$x' \in X_{k^{1/p}}$ the unique point lying over $x$, and
\item $\mathcal{O}_{X_{k^{perf}}, x'}$ is reduced for
$x' \in X_{k^{perf}}$ the unique point lying over $x$.
\end{enumerate}
\end{lemma}

\begin{proof}
Note that if $k'/k$ is purely inseparable, then
$X_{k'} \to X$ induces a homeomorphism on underlying topological
spaces, see Algebra, Lemma \ref{algebra-lemma-p-ring-map}.
Whence the uniqueness of $x'$ lying over $x$ mentioned in the
statement. Moreover, in this case
$\mathcal{O}_{X_{k'}, x'} = \mathcal{O}_{X, x} \otimes_k k'$.
Hence the lemma follows from Lemma \ref{lemma-geometrically-reduced-at-point}
above and Algebra, Lemma
\ref{algebra-lemma-geometrically-reduced-finite-purely-inseparable-extension}.
\end{proof}

\end{document}
